% SPDX-License-Identifier: Apache-2.0
\section{Method}
\label{sec:method}

\subsection{Minimal pairs}

The corpus is a set of small designs, each simulated once, each
differing from another case in exactly one property. A finding is the
difference between the two databases that results.

Two examples. \texttt{t10\_vec274} declares a value of 274 bytes and
\texttt{t10\_vec275} declares 275; the first is written as one record
and the second as two, which locates the threshold at which a large
value is split into chunks. \texttt{t2\_bit} and
\texttt{t1\_bit\_one\_edge} differ in the type of a single signal, and
the difference between the two type tables is the entry that
distinguishes \texttt{BIT} from \texttt{STD\_ULOGIC}.

Cases are grouped into tiers, each tier asking one question. The
corpus reached 1146 cases in 79 tiers. Each case directory holds the
source, a \texttt{BUILD.bazel} that simulates it, and a
\texttt{truth.json}.

\subsection{Ground truth declared before the file is opened}

\texttt{truth.json} states what the simulation must produce: the
signals with their scopes, widths and type names, the value of each at
each time, and, where a tier needs it, the number of raw records a
signal holds. It is derived from the design source. It is not derived
from the database, and it is written before the database is decoded.

This is what makes a decoding claim falsifiable. The reader must
reproduce the truth file; a decoding that merely reproduces itself
fails.

\subsection{The noise mask}

Not every byte of a database is a function of the design. Timestamps,
absolute paths, and several words that differ between two runs of the
same design would otherwise appear as differences between cases. A
mask records the fields known to vary, and the differential tooling
excludes them. Where a masked field was later explained, it left the
mask; where it was not, the report states it as unexplained rather
than as a constant.

\subsection{Recording rules}

Three rules govern what enters the format documentation.

\begin{enumerate}
  \item A fact enters the findings table only with the command that
        reproduces it, and the table names the case that found it and
        the cases that confirm it. Everything else stays in an open
        questions section and is labelled a guess.
  \item A test asserts against \texttt{truth.json}, or against
        Vivado's VCD as read by an existing parser. A test that
        asserts against bytes the decoder itself produced proves
        nothing and is not written.
  \item Every claim is scoped to the Vivado version that produced the
        file.
\end{enumerate}

A second table records which comparison produced which discovery, with
roughly 650 rows at the time of writing. That table is what makes the
result reviewable by a reader who was not present: a claim can be
traced to the pair of cases that produced it and re-run.

\subsection{Build environment}

Everything is built by Bazel: the Vivado installation is provisioned
by \texttt{rules\_vivado} into a shared cache, each corpus case is a
simulation target, and the test suite consumes the resulting databases
as data. A case is therefore reproducible from a clean checkout on any
machine with the same Vivado installer archive, and the whole corpus
is re-simulated by one command.
